% ============================================================================
% CHƯƠNG 10: KẾT LUẬN VÀ HƯỚNG PHÁT TRIỂN
% ============================================================================
\chapter{Kết luận}
\label{chap:conclusion}

\section{Tổng kết nội dung}

Báo cáo đã trình bày toàn diện về lập trình song song với CUDA và C\#, bao gồm các chủ đề:

\subsection{CUDA Programming}

\begin{enumerate}
    \item \textbf{Kiến trúc GPU}: Streaming Multiprocessors, CUDA cores, memory hierarchy
    \item \textbf{Mô hình lập trình}: Grid-Block-Thread, kernel functions
    \item \textbf{Các bài toán cơ bản}:
    \begin{itemize}
        \item Vector Sum và Matrix Sum (1D/2D indexing)
        \item Parallel Reduction (3 thuật toán với tối ưu khác nhau)
        \item Dot Product (multiply + reduce)
        \item Matrix Transpose (memory coalescing)
    \end{itemize}
    \item \textbf{Thuật toán sorting}:
    \begin{itemize}
        \item Histogram với atomic operations
        \item Counting Sort (stable, $O(n+k)$)
        \item Radix Sort (digit-by-digit)
        \item Quick Sort song song (partition + recursion)
    \end{itemize}
    \item \textbf{Graph algorithms}: BFS trong mạng đồng bộ với CSR format
\end{enumerate}

\subsection{C\# Parallel Programming}

\begin{enumerate}
    \item \textbf{Parallel.For}: Data parallelism với local state
    \item \textbf{Channel}: Producer-Consumer pattern
    \item \textbf{PLINQ}: Parallel LINQ queries
\end{enumerate}

\section{Kết quả đạt được}

\begin{table}[H]
\centering
\caption{Tổng hợp Speedup đạt được}
\label{tab:summary_speedup}
\begin{tabular}{@{}lccc@{}}
\toprule
\textbf{Bài toán} & \textbf{CPU (ms)} & \textbf{GPU (ms)} & \textbf{Speedup} \\
\midrule
Vector Sum (N=1M) & 2.5 & 0.08 & 31x \\
Matrix Sum (1024x1024) & 3.2 & 0.15 & 21x \\
Parallel Reduction (N=1M) & 1.8 & 0.05 & 36x \\
Matrix Transpose (4Kx4K) & 12.5 & 1.42 & 8.8x \\
Counting Sort (N=10M) & 450 & 18.5 & 24x \\
BFS (1M nodes) & 892 & 42.5 & 21x \\
\bottomrule
\end{tabular}
\end{table}

\begin{figure}[H]
\centering
\begin{tikzpicture}
\begin{axis}[
    width=13cm,
    height=8cm,
    xlabel={Bài toán},
    ylabel={Speedup (lần)},
    title={Tổng hợp Speedup GPU so với CPU},
    ybar,
    bar width=20pt,
    symbolic x coords={VectorSum, MatrixSum, Reduction, Transpose, CountSort, BFS},
    xtick=data,
    xticklabel style={rotate=20, anchor=east},
    nodes near coords,
    ymin=0,
    ymax=45,
    legend pos=north east,
]
\addplot[fill=blue!60] coordinates {
    (VectorSum, 31)
    (MatrixSum, 21)
    (Reduction, 36)
    (Transpose, 8.8)
    (CountSort, 24)
    (BFS, 21)
};
\addlegendentry{GPU Speedup}

\draw[dashed, red, thick] (axis cs:VectorSum, 20) -- (axis cs:BFS, 20);
\node[red, right] at (axis cs:BFS, 20) {Avg: 23.6x};

\end{axis}
\end{tikzpicture}
\caption{Biểu đồ tổng hợp speedup các bài toán}
\label{fig:summary_chart}
\end{figure}

\section{Bài học kinh nghiệm}

\subsection{Khi nào nên dùng GPU}

\begin{itemize}
    \item \textbf{Phù hợp}: Data parallelism, embarrassingly parallel, matrix/vector operations
    \item \textbf{Không phù hợp}: Sequential dependencies, small data, branch-heavy code
\end{itemize}

\subsection{Các kỹ thuật tối ưu quan trọng}

\begin{enumerate}
    \item \textbf{Memory Coalescing}: Truy cập bộ nhớ theo pattern liên tiếp
    \item \textbf{Shared Memory}: Cache data để giảm global memory access
    \item \textbf{Bank Conflicts}: Tránh multiple threads truy cập cùng bank
    \item \textbf{Warp Divergence}: Giảm thiểu branches trong warp
    \item \textbf{Occupancy}: Cân bằng giữa registers, shared memory và threads
\end{enumerate}

\begin{figure}[H]
\centering
\begin{tikzpicture}[scale=0.9]
    % Optimization checklist
    \node[font=\bfseries] at (5, 6.5) {CUDA Optimization Checklist};
    
    \draw[fill=green!10, rounded corners] (0, 0) rectangle (10, 6);
    
    \node[anchor=west] at (0.5, 5.5) {\checkmark \textbf{Memory Coalescing} - Threads truy cập địa chỉ liên tiếp};
    \node[anchor=west] at (0.5, 4.7) {\checkmark \textbf{Shared Memory} - Cache frequently accessed data};
    \node[anchor=west] at (0.5, 3.9) {\checkmark \textbf{Avoid Bank Conflicts} - Padding, stride access};
    \node[anchor=west] at (0.5, 3.1) {\checkmark \textbf{Minimize Warp Divergence} - Uniform control flow};
    \node[anchor=west] at (0.5, 2.3) {\checkmark \textbf{Optimize Block Size} - Thường 128-256 threads};
    \node[anchor=west] at (0.5, 1.5) {\checkmark \textbf{Reduce Memory Transfers} - Compute on GPU};
    \node[anchor=west] at (0.5, 0.7) {\checkmark \textbf{Use Streams} - Overlap compute và transfer};
    
\end{tikzpicture}
\caption{Checklist tối ưu CUDA}
\label{fig:optimization_checklist}
\end{figure}

\section{Hướng phát triển}

\subsection{Mở rộng kiến thức}

\begin{enumerate}
    \item \textbf{Advanced CUDA}:
    \begin{itemize}
        \item CUDA Streams và Concurrent Execution
        \item Multi-GPU programming
        \item Unified Memory
        \item Dynamic Parallelism
    \end{itemize}
    
    \item \textbf{Libraries}:
    \begin{itemize}
        \item cuBLAS - Linear algebra
        \item cuFFT - Fast Fourier Transform
        \item Thrust - High-level algorithms
        \item cuDNN - Deep learning
    \end{itemize}
    
    \item \textbf{Ứng dụng thực tế}:
    \begin{itemize}
        \item Machine Learning / Deep Learning
        \item Image Processing
        \item Scientific Simulation
        \item Cryptography
    \end{itemize}
\end{enumerate}

\subsection{So sánh với các công nghệ khác}

\begin{table}[H]
\centering
\caption{So sánh các nền tảng lập trình song song}
\label{tab:platforms_comparison}
\begin{tabular}{@{}lcccc@{}}
\toprule
\textbf{Tiêu chí} & \textbf{CUDA} & \textbf{OpenCL} & \textbf{OpenMP} & \textbf{MPI} \\
\midrule
Platform & NVIDIA GPU & Any GPU/CPU & CPU (multi-core) & Distributed \\
Ease of use & Medium & Hard & Easy & Medium \\
Performance & Excellent & Good & Good & Variable \\
Portability & NVIDIA only & Cross-platform & Cross-platform & Cross-platform \\
Use case & GPU compute & Heterogeneous & Shared memory & Clusters \\
\bottomrule
\end{tabular}
\end{table}

\section{Kết luận}

Lập trình song song là kỹ năng thiết yếu trong thời đại big data và AI. Qua môn học CT225H và báo cáo này, sinh viên đã:

\begin{enumerate}
    \item Nắm vững \textbf{kiến trúc GPU} và mô hình lập trình CUDA
    \item Thực hành các \textbf{thuật toán song song cơ bản} đến nâng cao
    \item Hiểu các \textbf{kỹ thuật tối ưu} để đạt hiệu năng tối đa
    \item So sánh \textbf{CPU vs GPU} và biết khi nào nên sử dụng mỗi loại
    \item Ứng dụng \textbf{C\# Parallel} cho các bài toán trên CPU
\end{enumerate}

\vspace{1cm}

\begin{center}
\textit{``The future of computing is parallel, and the future of parallel computing is on GPU.''}

\vspace{0.5cm}
--- NVIDIA
\end{center}

\vspace{1cm}

\noindent\textbf{Lời cảm ơn:}

Em xin chân thành cảm ơn \textbf{ThS. Lâm Hoài Bảo} đã tận tình hướng dẫn và truyền đạt kiến thức quý báu trong suốt môn học CT225H - Lập trình Song song. Những kiến thức này sẽ là nền tảng vững chắc cho con đường nghiên cứu và phát triển trong tương lai.
